\section{Material und Methoden}
\textcolor{red}{• Beschreibung des Datensatzes • statistisches Modell!!!
• deterministischer und stochastischer Teil
• verwendete statistische Methoden
• Beschreibung der Vorgehensweise
• kein Beweis theoretischer Resultate, aber Definition der verwendeten Größen
• nachvollziehbare Beschreibung der Simulation}

\subsection{Datensatz}

Die Analyse basiert auf dem Datensatz \texttt{hainich.csv}, der von A.~Soe (MPI Biogeochemie)
bereitgestellt wurde. Der Datensatz enthält 38 Messungen der Bodenatmung (\texttt{soil.res}) sowie
folgende Einflussgrößen:

\begin{itemize}
\setlength{\itemsep}{0pt}
\item \texttt{temp.0}, \texttt{temp.5}, \texttt{temp.10}, \texttt{temp.15} – Bodentemperatur in 0, 5, 10, 15 cm Tiefe
\item \texttt{smoi} – relative Bodenfeuchte (soil moisture)
\item \texttt{lmoi} – relative Feuchte der Streuschicht (litter moisture)
\item \texttt{litter.d} – Dicke der Streuschicht
\item \texttt{ph.litter}, \texttt{ph.soil0}, \texttt{ph.soil5} – pH-Wert der Streuschicht bzw. des Bodens in 0-5 cm bzw. 5-10 cm Tiefe
\item \texttt{rootdw0}, \texttt{rootdw5} – Trockenmasse der Wurzeln in 0-5 cm bzw. 5-10 cm Tiefe
\item \texttt{rootn0}, \texttt{rootn5} – Stickstoffgehalt der Wurzeln in 0-5 cm bzw. 5-10 cm Tiefe
\item \texttt{rootc0}, \texttt{rootc5} – Kohlenstoffgehalt der Wurzeln in 0-5 cm bzw. 5-10 cm Tiefe
\item \texttt{litter.c}, \texttt{litter.n} – Kohlenstoff- bzw. Stickstoffgehalt der Streuschicht
\item \texttt{soiln0}, \texttt{soiln5} – Stickstoffgehalt des Bodens in 0-5 cm bzw. 5-10 cm Tiefe
\item \texttt{soilc0}, \texttt{soilc5} – Kohlenstoffgehalt der Bodens in 0-5 cm bzw. 5-10 cm Tiefe
\item \texttt{litno3}, \texttt{sno30}, \texttt{sno35} – Nitratgehalt der Streuschicht bzw. des Bodens in 0-5 cm bzw. 5-10 cm Tiefe
\item \texttt{litnh4}, \texttt{snh40}, \texttt{snh45} - Ammoniumgehalt der Streuschicht bzw. des Bodens in 0-5 cm bzw. 5-10 cm Tiefe
\item \texttt{litdoc}, \texttt{sdoc0}, \texttt{sdoc5} – Gehalt an gelöstem organischen Kohlenstoff (DOC - dissolved organic carbon) der Streuschicht bzw. des Bodens in 0-5 cm bzw. 5-10 cm Tiefe
\item \texttt{nn} – Entfernung der Probestelle zum nächsten Baum (nearest neighbour)
\end{itemize}

\subsubsection{Transformation}
Die Zielgröße \texttt{soil.res} zeigt eine linksschiefe Verteilung. Um eine möglichst stabile Schätzung des Modells
zu erreichen, wurden die Zielgröße mit dem natürlichen Logarithmus transformiert.

\subsubsection{Interaktionen zwischen den Einflussgrößen} \label{Interaktionen}

Es gibt Einflussgrößen deren Effekt, den sie auf die Zielgröße ausüben, von der Ausprägung
einer anderen Einflussgröße abhängig ist. Dabei handelt es sich um eine Interaktion, welche in einem
linearem Modell explizit beschrieben werden muss.

folgende Interaktionen zwischen den Einflussgrößen gelten als biologisch relevant:
\\
\texttt{lmoi} und \texttt{litter.de}:\\
Je dicker die Streuschicht ist, desto mehr Wasser kann sie speichern und desto länger dauert es,
bis sie ausgetrocknet ist.
\\
\texttt{temp.15} und \texttt{lmoi}:\\
Mikroorganismen drosseln bei extremer Trockenheit trotz hoher Temperaturen ihre Aktivität
In trockenen Perioden (wie im August 2001 beobachtet) bricht die normale
Temperatur-Respirations-Kurve ein. Das Modell für Bodenatmung wird signifikant besser,
wenn sowohl Temperatur als auch Feuchte einbezogen werden
\\
\texttt{temp.15} und \texttt{rootdw0}:\\
Die Temperatur-Sensitivität (Q10) unterscheidet sich zwischen der autotrophen Respiration (Wurzeln) und der heterotrophen Respiration (Mikroorganismen)
Wurzelrespiration macht oft mehr als die Hälfte der gesamten Bodenatmung aus. Da Wurzeln und
Mikroorganismen unterschiedliche physiologische Optima haben, hängt die Gesamtreaktion des Bodens auf
Erwärmung davon ab, wie viel Wurzelbiomasse vorhanden ist
\\
\texttt{soilc0} und \texttt{rootdw0}:\\
Räumliche Muster der Bodenatmung werden oft gemeinsam durch die Verteilung von Pflanzen (Wurzeln) und Bodenressourcen (C) erklärt
Das beste Modell in der Publikation kombiniert explizit Wurzelbiomasse, Boden-C und Feuchte.
Diese Faktoren korrelieren oft räumlich miteinander, da Bäume über den Wurzeleintrag den Bodenkohlenstoff beeinflussen
\\
\texttt{litter.d} und \texttt{soilc0}:\\
Die Verfügbarkeit von Substrat (Kohlenstoff) für den Abbau wird durch die Feuchtigkeit limitiert oder gefördert.
Ein sehr hoher Wassergehalt (besonders in tonreichen Böden wie im Hainich) kann die
Sauerstoffdiffusion hemmen und so die mikrobielle Aktivität trotz vorhandenen Kohlenstoffs
einschränken. Zudem sind Streufeuchte (lmoi) und Streudicke (litter.d) zentrale Parameter für die
räumliche Variabilität

\subsection{Statistisches Modell} \label{Statistisches Modell}

Zur Vorhersage der Bodenatmung wurde ein lineares Modell (\texttt{lm}) verwendet.\\
Deterministischer Teil: $\mathbb{E}Y_i= \mu_i =\beta_1 + \beta_2x_{i2} + ... + \beta_kx_{ik}$ mit $i= 1,...,n$\\
Stochastischer Teil: $Y_i \sim N(\mu_i, \sigma^2)$ mit $i= 1,...,n$\\
\\ 
Um die Anzahl der potenziellen Prädiktoren auf ein für die Größe des Datensatzes angemessenes Maß zu
reduzieren, wurde die gängige Heuristik angewandt, nur Variablen mit einer Einzelkorrelation zum Betrag
von >0,2 für die weitere Modellierung zu berücksichtigen. Dies ermöglichte eine Konzentration auf die
einflussreichsten Faktoren und diente dem Schutz vor Overfitting.
Die globale Suche wurde mit der Funktion \texttt{regsubsets()} durchgeführt. Diese Funktion testet alle möglichen Kombinationen der
Einflussgrößen in einem linearem Modell.
\\
Weiterhin wurden alle biologisch plausiblen Interaktionen getestet, siehe \ref{Interaktionen}. Die Interaktion die den AIC verbessert hat, wurde
in das Modell einbezogen. Durch die Menge der Daten, ist die Interaktion nicht statistisch signifikant. Wurde jedoch trotzdem behalten aufgrund des verbessertem AIC und dem Parsimonie-Prinzip.

\subsubsection{Validierung des finalen Modells}

Zur Beurteilung der Qualität der Modelle wurden mehrere Validierungsschritte durchgeführt. Ziel war es, sowohl die statistische Zulässigkeit der Modelle als auch deren Prognosegüte und Komplexität zu bewerten.\\

\subsubsection{Überprüfung der Modellannahmen}
Die Standard-Diagnoseplots des linearen Modells wurden betrachtet (Residuals vs. Fitted, Normal Q-Q, Scale-Location, Residuals vs. Leverage). Diese dienen der Identifikation nichtlinearer Strukturen, der Beurteilung der Varianzhomogenität sowie der Erkennung einflussreicher Beobachtungen. \\

Zur formalen Prüfung der Normalverteilung der Residuen wurde der Shapiro-Wilk-Test verwendet. Die Teststatistik ergibt sich zu
\[W = \frac{\left( \sum_{i=1}^n a_i \varepsilon_{(i)} \right)^2}{\sum_{i=1}^n (\varepsilon_i - \bar{\varepsilon})^2},\]
wobei \(\varepsilon_{(i)}\) die geordneten Residuen und \(a_i\) konstante Gewichte sind. Werte von \(W\) nahe \(1\) sprechen für eine Normalverteilung.\\

Zur Prüfung der Homoskedastizität wurde der Breusch-Pagan-Test verwendet, welcher in R mit der Funktion \texttt{bptest()} implementiert ist. Dieser untersucht, ob die Varianz der Residuen von den Regressoren abhängt. Hierzu wird eine Hilfsregression der quadrierten Residuen auf die erklärenden Variablen durchgeführt:
\[\hat{u}_i^2 = \gamma_0 + \gamma_1 z_{i1} + \dots + \gamma_q z_{iq} + v_i.\]
Unter der Nullhypothese \(H_0: \gamma_1 = \dots = \gamma_q = 0\) liegt Homoskedastizität vor, während signifikante Koeffizienten auf Heteroskedastizität hinweisen.\\

Um Einfluss einzelner Beobachtungen auf die Parameterschätzung zu messen, wurde die Cook's Distance verwendet, berechnet in R mit \texttt{cooks.distance()}:
$$D_i = \frac{\sum_{j=1}^n (\hat{y}_j - \hat{y}_{j(i)})^2}{(k + 1) s^2}.$$
Große Werte von \(D_i\) weisen auf Beobachtungen hin, die die Modellschätzung stark beeinflussen.\\

\subsubsection{Modellwahl}
Für die Qualität der Modelle wurde das BIC genutzt. Das BIC ist das Bayesian Information Criterion und ist
ein Kriterium zur Modellwahl. Es wägt die Anpassungsgüte eines statistischen Modells gegen seine Komplexität, also die Anzahl der
Variablen ab. Das BIC wird berechnet mit der Formel: \[-2\cdot logL+log(n)\cdot p \] mit der Log-Likelihood (L) des Modells, der Anzahl der
geschätzten Parameter (p) und dem Stichprobenumfang (n). Damit lässt sich einfach der Varianz Trade-off ablesen.\\

\subsubsection{Bewertung der Prognosegüte}
Zur Bewertung der Prognosegüte wurden LOOCV (Leave-One-Out-Cross-Validation) sowie k-fold-Cross-Validation eingesetzt.\\

Bei der LOOCV wird das Modell \(n\)-mal geschätzt, wobei jeweils eine Beobachtung als Testfall und die übrigen \(n-1\) als Trainingsdaten verwendet werden. Dadurch wird jede Beobachtung genau einmal validiert und der Vorhersagefehler auf unabhängigen Daten geschätzt.

\[\mathrm{RMSE}_{\mathrm{LOOCV}} = \sqrt{\frac{1}{n} \sum_{i=1}^n (y_i - \hat{y}_{(i)})^2}\] \\

Bei der \( k \)-fold-Cross-Validation wird der Datensatz in \( k \) disjunkte Teilmengen (Folds) aufgeteilt. In jeder Iteration dient ein Fold als Testmenge, während die übrigen \( k-1 \) Folds zur Modellschätzung verwendet werden. Dieser Vorgang wird \( k \)-mal wiederholt. Für jeden Durchlauf wird ein mittlerer quadratischer Fehler \(MSE_j\) bestimmt, der auf dem jeweiligen Testfold basiert.

\[\mathrm{RMSE}_{k\text{-fold}} = \sqrt{\frac{1}{k} \sum_{j=1}^{k} MSE_j}\]\\

Durch die Kombination beider Verfahren kann sowohl eine nahezu unverzerrte als auch eine stabilere Schätzung des Vorhersagefehlers erzielt werden.

Gegenstand der Simulationsaufgabe ist die falsche Anwendung des t-Tests bei der
Vorwärtsselektion von Prädiktoren. Während der Vorwärtsselektion wird mehrfach
eine Hypothese der Form
überprüft, wobei jedoch jeweils die absolut größte von mehreren möglichen t-
Statistiken betrachtet wird.

Wählen Sie das im ersten Teil bestimmte Modell als wahres Modell für die
Simulation. Wählen Sie für jeden Stichprobenumfang eine feste Designmatrix , die aus
zufällig ausgewählten Zeilen der ursprünglichen Designmatrix besteht (evtl. mit
Wiederholung). Simulieren Sie dann jeweils mehrfach Pseudobeobachtungen der
Zielgröße. Vergleichen Sie die simulierte Häufigkeitsverteilung des Maximums
mehrerer t-Statistiken für die Hinzunahme einer weiteren Einflussgröße beim nächsten
Schritt der Vorwärtsselektion mit der für einen einzelnen Test passenden t-Verteilung
und diskutieren Sie die Auswirkungen auf die Wahrscheinlichkeit einer irrtümlichen
Ablehnung der Null-Hypothese.


\subsection{Simulationsstudie}

Ziel der Simulationsstudie ist eine mehrfache Vorwärtsselektion von Prädiktoren.

Nach der wahl des besten Modells ($\beta$) in \ref{Statistisches Modell} erfolgt die berechnung einer Designmatrix ($X$). Diese wird fuer jede Simulation neu gezogen ist aber für jeden Stichprobenumfang in einer Simulation gleich. Die berechnung geschieht durch ziehen mit zurücklegen aus der ursprünglichen Designmatrix. 
Die simulierung mehrerer Pseudobeobachtungen ($Y$) ist notwendig um, die verteilung von Teststatistiken zu approximieren. Simuliert werden diese durch:\\
\begin{equation}
\mathrm{Y} \sim X\beta + \epsilon
\end{equation}
wobei $\epsilon$ ein gleichverteilte zufallszahl aus $[-1,1]$ ist. 

Danach folgt eine serie an T-tests wobei jeder T-test die qualität der aufnahme eines praediktors in das modell tested. Dies passiert durch testung der Null Hypothese:\\
\begin{equation}
\mathrm{H_0 : \beta_j}  = 0 
\end{equation}
was zusammen dem Einstichprobentest: 
\begin{equation}
t = \frac{\overline{x}-\mu_0}{\frac{s}{\sqrt{n}}}
\end{equation}
%\[\mathrm{t_{\hat{\beta}} =} \frac{\hat{\beta} - \beta_0}{\mathrm{SE}(\hat{\beta})}\]
den T-statistik bildet:
\begin{equation}
T = \mathop{max}_{\textbf{j}}|t_j|
\end{equation}

Wir selektieren den praediktor mit dem hoechesten tmax und testen ob der T wert ueber $\alpha$ liegt. Wenn dies der fall ist nehmen wir den praediktor in unser modell auf und beginnen eine neues selektionsverfahren. Wenn nicht behalten wir das vorherige modell und versuchen die aufnahme erneut. Dies passiert bis zu einer vordefinierten anzahl an aufnahme schritten. 

\subsubsection{Parameter und Parameterwahl}
Die Simulationsstudie ist stark abhaengig von der wahl der parameter. Die Parameter der Simulation die gesetzt werden koennen sind:
\begin{itemize}
    \item Stichprobenumfang
    \item Anzahl der Simulationen
    \item Fehlerstandardabweichung $\sigma$
    \item Anzahl der Selektionsschritte
    \item Kritischer wert
\end{itemize}

Der Stichprobenumfang beschreibt wie haeufig wir aus der ursprünglichen Designmatrix ziehen um die Designmatrix fuer jede Simulation zu erstellen. Dafuer haben wir den wert 100 gewaehlt da dies zu stabilieren schaetzungen fuehrt als kleinere Stichprobenumfaenge wie 20 oder die groese der originalen designmatrix (38). In formel \ref{formel 3} ist der Stichprobenumfang im zaehler des zaehlers, es folgt das die erwarteten T-werte groeser sein werden wenn der Stichprobenumfang steigt. 

Die anzahl an Simulationen beschreibt wie viele simulationen in betracht gezogen werden bei der Wahl des naechsten prediktor der aufgenommen wird. Diesen wert haben wir auf 50,000 gesetzt. Hier gilt grundsaetzlich groeser ist besser allerdings hat mehr als 50,000 die laufzeit zu sehr angehoben also mussten wir da aufhoeren.

Die Fehlerstandardabweichung haben wir auf 1 gesetzts was gaengige praxis entspricht das gleiche gilt fuer den kritischen wert von 0.05.

Die anzahl an selektionsschritten haben wir auf 10 gesetzt da im schnitt 7 praediktoren aufgenommen werden und wir noch ein paar schritte simulieren wollten um zu zeigen das es nicht nur zufall war das es nichts mehr aufgenommen hat.

\subsection{Code-Verfügbarkeit}
link zum repository\\
TODO MAKE REPO PUBLIC